% Generated from validated Article Draft Result. Do not hand-edit; revise the structured draft instead.
\begingroup\raggedright
\subsection{モデル更新を「性能表」ではなく利用面で読む}
\par\endgroup

3月のfrontier modelは、単に新しい名前が増えた月ではない。OpenAIではGPT-5.3 InstantからGPT-5.4、さらにmini/nanoへ月内でfamilyが広がり、Anthropic、Google、Z.ai、Alibabaでもmodelやdeveloper platformの更新が確認できる。比較軸は単一benchmarkではなく、tool use、computer use、multimodal、API lifecycle、提供形態へ広がっていた。\autocite{src-46b3e364b737,src-8afdb503cc56,src-9156c37fb1c2,src-a24fd97eb3d0,src-eac8dfbeb78c,src-eea6230f612a,src-f6aa679babd7}


\begingroup\raggedright
\subsection{OpenAI: GPT-5.3 InstantからGPT-5.4 familyへ}
\par\endgroup

一次資料では3月3日にGPT-5.3 Instant、3月5日にGPT-5.4、3月17日にGPT-5.4 mini/nanoが並ぶ。OpenAIはGPT-5.4をcoding、computer use、tool search、長いcontextなどに位置付け、mini/nanoにはより小型・高速な用途やsubagent利用を挙げた。ここで確実なのは月内chronologyとvendorが示したscopeであり、性能優位そのものではない。\autocite{src-46b3e364b737,src-8afdb503cc56,src-eea6230f612a}


\begingroup\raggedright
\subsection{Claude、Gemini、GLM、Alibaba: 更新面は一様ではない}
\par\endgroup

Anthropicの保存済みnews indexには3月31日のClaude Sonnet 4.6が記録される。Gemini APIでは3月3日から31日にかけて軽量model、multimodal embedding、built-in toolsとfunction calling、live audio、video生成系まで複数回の更新が並んだ。Z.aiは3月15日のGLM-5-Turboをlong-chain agent workloadやtool連携に位置付け、Alibaba Model Studioのsnapshotには3月3日のQwen Image 2.0系などが残る。3月はmodel identityだけでなく、modelをどのsurfaceで使うかが同時に更新された月だった。\autocite{src-9156c37fb1c2,src-a24fd97eb3d0,src-eac8dfbeb78c,src-f6aa679babd7}


したがって3月のfrontier領域は、モデル単体の能力競争を続けながら、その能力をtool、computer environment、multimodal surface、API lifecycleへどう接続するかという競争へ広がっていた、と読むのが妥当だ。この見方は各社のbenchmarkを横断ランキングへ変換するものではない。\autocite{src-a24fd97eb3d0,src-eea6230f612a,src-f6aa679babd7}


